% =============================================================================
% 5 Implementation. Trimmed 2026-10-19 to OSDI/MLSys norms (Ben): two
% paragraphs, ~0.4-0.5 page. Only what s4 does not already say and a reproducer
% would trip over. No implementation gaps in the body (Appendix C). The handover
% emulation is described in the s7 Setup; the ablation switches are pointed to
% from the s7 Method.
% =============================================================================

\section{Implementation}\label{sec:impl}

\paragraph{Prototype.} \sysname{} is a fork of vLLM's V1 engine at upstream
commit \texttt{b89275d01}. The rider path is about 5{,}800 lines in nine
files. It touches the scheduler, the model runner and the FlashAttention
backend, and adds a host-pool module; no model code is modified. The prototype
runs on a single GPU, and the design is per GPU: under tensor parallelism each
GPU holds $1/N$ of the weights and of every request's KV and has its own link,
so the capacity model applies to each GPU separately. The engine runs in eager
mode, without CUDA-graph capture, and every baseline in this paper runs in the
same eager mode.
% src: git diff --numstat b89275d01..HEAD (upstream base 2026-01-15):
%      hybrid_rungd_pool.py 1608, gpu_model_runner.py +789, scheduler.py +519,
%      flash_attn.py +478, transfer_aware_admission.py 949, kv_tier_backend.py
%      248, hybrid_attn_ctx.py 100, hybrid_pool_admit.py 94,
%      pipelined_decode_bench_connector.py 1064 -> 5849 added.
%      hybrid_attn_ctx.py docstring: "a KV connector forces eager execution (no
%      CUDA-graph capture)". Baselines: r=0 rows of rider_cost_curve.py, same
%      engine and flags.

\paragraph{Engineering details.} The host copy of every rider's KV is pinned
memory, allocated once for all $L$ layers and viewed per layer. Pinning is what
lets each run be a single \texttt{cudaMemcpyAsync} straight into the slot: a
copy from pageable memory is staged through a driver buffer and blocks the
host. A single allocation lets the deferred scatter
(\cref{sec:rolling}) write every layer's new rows with two vectorised index
copies, one for $K$ and one for $V$, rather than $2L$, because a rider's row
in its extent is the same in every layer. The write-back needs no stream of
its own. Its device-to-host copies are issued non-blocking on the compute
stream behind one recorded event, and only the host-side scatter waits on that
event at the start of the next step. The second stream carries the
host-to-device prefetch alone.
% src: hybrid_rungd_pool.py: RungdHostPool.__init__ (_unified_mirror:
%      _host_k_all/_host_v_all pin_memory=True, per-layer views; _wt_layer_off
%      "global row = host_pos + layer*h_max*bs"; _copy_stream is the only
%      torch.cuda.Stream), write_token (copy_(non_blocking=True) on the current
%      stream, then self._wt_event.record()), flush_token_writes
%      (self._wt_event.synchronize() before the scatter); the code comment on
%      the pre-T2 path: "two PAGEABLE (hence blocking) .to("cpu") copies".
